\section{Construction}\label{sec:offline-construction}

In brief, the protocol proceeds as follows. Each user enrols once with the
registration authority, which binds a secret key held in their secure element to
a unique identifier. A user withdraws a token through their intermediary, and
the central bank signs it without learning what it has signed. A payment
transfers ownership of the token from payer to payee while both are
disconnected, and extends a chain of anonymous signatures that the token carries
with it. The final owner deposits the token through their intermediary, and the
central bank detects any token that is deposited twice.  An auditor then
recovers the identity of the party responsible for the duplicate, and of no one
else.

The construction assumes that the registration authority issues each user
exactly one identifier, and issues a credential only to a secure element that
holds the corresponding secret key. This is what prevents a user from
accumulating identities, and it is what makes identification meaningful: an
identifier recovered at audit names one enrolled user. We assume identifiers are
unpredictable and that a user does not publish their own. We assume the secure
element does not disclose its secret key. We do not assume it succeeds in
enforcing that a token is spent once. That enforcement is what makes double
spending expensive to attempt, but detection and identification do not rest on
it, and a scheme that relied on hardware alone would have no recourse once the
hardware failed. We assume payees hold the public keys of the central bank and
of the registration authority, so that they can verify a payment without
contacting either. Addressing a malicious registration authority is orthogonal
to our contribution. We now present the construction in full.

\paragraph{Tokens.} A token is a value together with a pseudonym for its owner,
and it is transferred whole. It is never split and never merged, which is what
keeps a payment a single fixed size operation rather than one that scales with
the amount being paid.

The pseudonym is a hiding commitment to the owner's identifier. Committing
rather than naming a public key is what makes a payment anonymous to anyone who
reads its transcript, since a user commits afresh for every token they hold.
The commitment randomness does a second job. It is what distinguishes two
payments that are otherwise identical, so it never has to be sent alongside the
pseudonym.

\paragraph{Setup.} During setup the central bank generates a key pair for a
blind signature scheme, the registration authority generates a key pair for a
signature scheme over message vectors, and the auditor generates an encryption
key pair. All public keys are published.

Both signature schemes sign vectors, because a token binds an owner, a value,
and the commitment randomness, and a credential binds an identifier to a secret
key. The auditor's key is only ever used to decrypt, and only after a conflict
has already surfaced. That is what allows the role to be distributed across
several parties by threshold decryption, without any change to the rest of the
protocol.

The generators that make up the commitment key are derived by hashing into the
group rather than chosen by any participant, so that no party knows the discrete
logarithm relations between them. Binding of the commitments is what the
identification argument rests on, and a central bank that generated the
commitment key itself would be able to break it.

\paragraph{Registration.} A user enrols once. Their secure element is
provisioned with a fresh secret key and a unique identifier, and proves to the
registration authority that it holds that key. The registration authority issues
a credential binding the two, blindly, so that the secret key never leaves the
secure element even during enrolment. The user's mobile device is then paired
with the secure element.

The two checks do different work. Insisting the identifier is fresh is what
makes an identifier recovered at audit name one enrolled user rather than a role
several people share. Insisting the secure element proves knowledge of the key
is what binds the credential to hardware the enrolling user actually holds, so
that a payment authorised under that credential can be attributed to them.
Pairing the device to the secure element matters for a reason that only becomes
visible at payment time. The two can communicate directly, so the secure element
can pass the device values it would otherwise have to derive from a shared key.

\paragraph{Withdrawal.} To withdraw, the user's secure element generates a
pseudonym for the new token, and the device asks the intermediary for a
signature on it from the central bank. The intermediary checks the user has the
funds and forwards the request. The central bank checks the intermediary has the
reserves and signs. The device then converts the signature it receives into a
proof that the token was issued.

Two constraints shape this exchange. The central bank must be satisfied that the
token is well formed and that the funds exist, but it must not learn who the
token belongs to, or it could later match a withdrawal to a deposit. The request
therefore discloses the value and hides the identifier and the randomness, and
the signature is produced blindly. The value is disclosed because the central
bank is accounting for it. This is also the one channel that remains: a
withdrawal and a deposit of unusual value can still be matched, which is why a
deployment should use a small set of fixed denominations.

The second constraint is that the signature can never be shown. Verifying it
requires the messages it was issued on, and one of those is the withdrawer's
identifier, so a payee who checked the signature directly would learn who
withdrew the token. The device therefore proves knowledge of it instead, once,
at withdrawal, and the resulting proof of withdrawal is what every subsequent
payee verifies.

Nothing in the withdrawal binds the committed identifier to the withdrawing
user, since the central bank never sees it. It does not need to. Spending the
token requires a credential on that identifier, so a user who withdrew a token
committing to somebody else's identifier could never spend it, and would only
have destroyed their own funds.

\paragraph{Payment.} A payment begins with the payee. Their secure element
generates a fresh pseudonym and sends it to the payer. The payer's secure
element locks the token being spent, the payer's device encrypts the payer's
identifier under the auditor's key, and the two jointly produce a signature of
knowledge, over the value and the two pseudonyms, transferring the token to the
payee. The payer sends the payee the token, the ciphertext, the signature, and
everything the payee needs to check the token's history. On acknowledgement the
payer's secure element deletes the token.

Signing the payee's pseudonym along with the payer's is what fixes the
destination of the payment. Were only the payer's pseudonym signed, the
signature would carry no commitment to where the token was going, and anyone who
saw it could redirect the token to a pseudonym of their own.

Having the payee choose the new pseudonym is what makes identification safe.  If
the payer chose it, a payer who spent the same token to the same payee twice
could produce two identical payments, and the payee would then be
indistinguishable from the party that double spent. A payee cannot defend
themselves against this by remembering every payment they have received, since
that is not something a wallet can be expected to do. Because the payee
contributes randomness instead, the two payments differ, and the point at which
two histories diverge names the payer who sent them.

The payee draws that randomness uniformly, so two payments of the same token
collide only with negligible probability. A payee who repeats it deliberately
gains nothing. Both payments then agree at their pseudonym, the divergence moves
one link further along, and it is the payee who is identified as having spent
one token twice.

The ciphertext is present because a payee is offline and cannot ask anyone who
the payer is. Anything the auditor will need must already be in the transcript,
sealed. The payer's device therefore encrypts its own identifier and proves that
the ciphertext matches the pseudonym it is spending under.  That proof is what
stops a double spender sealing someone else's identity in place of their own. It
is also what protects an honest payee from a malicious payer who chooses their
pseudonym for them, since producing it requires knowing the identifier the
pseudonym commits to.

That proof and the transfer signature both open the payer's pseudonym, on the
same witness, and both are produced by the device. We therefore compile them as
a single proof over a shared challenge, so the opening is proved once per hop
rather than twice. We keep them as separate symbols below, because the transfer
signature is produced jointly with the secure element and the audit proof is
not, and merging them is a property of the instantiation rather than of the
construction.

The signature is produced jointly because the secret key never leaves the secure
element, and the secure element is not powerful enough to produce the whole
proof. The secure element contributes only knowledge of the key. The device
supplies the rest of the witness and does the work. This split is what makes the
scheme deployable on hardware that exists today, and it is developed
in~\Cref{sec:joint-proof}.

The payee verifies every link in the chain before accepting, since nobody else
will do it for them. This is the cost of operating offline: the transcript grows
by one link per transfer, and verification grows with it.
\Cref{sec:offline-evaluation} measures both.

\paragraph{Deposit.} The final owner signs the token over to the central bank
through their intermediary. The central bank verifies the chain, checks whether
it has seen a deposit of the same token before, and credits the account if it
has not.

A deposit carries the depositor's own audit ciphertext and proof, exactly as a
payment does, and its signature of knowledge is taken over a fixed deposit
marker in place of a payee pseudonym. Carrying the ciphertext is what lets the
auditor name a user who both deposits a token and pays it onward, or who
deposits it twice, where the two chains agree everywhere and one simply stops
earlier than the other. Using a marker distinct from any pseudonym is what stops
a deposit signature being replayed as a payment.

The check is against the token's original pseudonym, the one fixed at
withdrawal. That is the right thing to compare because it is the only part of a
token that transfers do not change. Two deposits of the same token agree there
whatever they do afterwards, so a conflict cannot be hidden by paying the token
along a different route.

A rejected deposit is not a loss borne by the party that presented it. The
central bank holds the conflicting requests, the auditor names the user at the
divergence, and the double spent value is settled against that user's account by
the central bank once they are named. Whoever presented the second deposit is
credited on the strength of a transcript that verifies, so an honest holder who
was paid a token that had already been spent elsewhere is made whole, and the
shortfall is carried by the party that created it. This differs from schemes in
which the issuer absorbs the double spent value, since here the total in the
system is unchanged: the credit to the depositor is matched by a debit to the
identified spender rather than by fresh issuance. The settlement itself is an
operation of the central bank rather than a step of the protocol, and further
consequences for the named user follow outside it.

\paragraph{Identification.} Two conflicting deposits of one token agree on a
prefix and then part company, either because the owner at that point sent two
different payments, or because they deposited the token and also paid it onward.
The auditor takes the last position at which the two agree, decrypts the
ciphertext the owner attached there, and recovers their identifier.

Nothing else is recovered. The ciphertexts before that position belong to owners
whose payments agree in both deposits, so each of them passed the token on
exactly once, and their ciphertexts are never opened. The auditor can in
principle deanonymise any payment, which is why we recommend distributing the
role so that decryption requires more than a threshold of auditors to cooperate.

\paragraph{Privacy guarantees.} The scheme gives three guarantees. Transactors
are anonymous to anyone who sees a transaction transcript, assuming payer and
payee exchange identities out of band. Withdrawals and deposits are unlinkable
to every party in the system, so no one can trace a deposit or a payment to its
origin or its destination from the transcripts. The auditor may deanonymise
payments, and is trusted to do so only on evidence of double spending.

The complete protocol is given in Protocol~\ref{prot:offline}. It is stated
generically with respect to the underlying commitment scheme, digital signature
scheme, blind signature scheme, public key encryption scheme, and joint proof of
knowledge. \Cref{sec:offline-instantiation} instantiates each of them.

\begin{protocolbox}{Offline payment protocol}{offline}

  \begin{description}

    \item[$(\pk_{\algo{cb}}, \ck, \pk_\auditor, \rapk) \leftarrow
\OPgen(\secparam)$ :] On input security parameter $\secparam$, do:

      \begin{enumerate}

	\item The central bank generates keys for a blind signature scheme
supporting three-message vectors, obtaining $(\sk_{\algo{cb}},
\pk_{\algo{cb}})$.

	\item Run $\ck \leftarrow \ComGen(\secparam, 2)$ to obtain a public
commitment key, with all generators derived by hashing into the group, so that
no participant knows the discrete logarithm relations between them.

	\item The auditor generates an asymmetric encryption key pair
$(\sk_\auditor, \pk_\auditor) \leftarrow \PKEgen(\secparam)$.

	\item The registration authority generates keys for a multi-message
signature scheme supporting two-message vectors, obtaining $(\rask, \rapk)
\leftarrow \DSgen(\secparam)$.

	\item Publish all public keys and generators to every participant, and
return $(\pk_{\algo{cb}}, \ck, \pk_\auditor, \rapk)$.

      \end{enumerate}

    \item[$\cred \leftarrow \OPregister(\rask)$ :] When a user opens an account,
do the following:

      \begin{enumerate}

	\item Provision the user's secure element $\algo{se}$ with a fresh
secret key $\sk_{\algo{usr}}$ and a unique random identifier $\algo{uid}$.

	\item $\algo{se}$ sends the registration authority $\algo{uid}$ and a
commitment to $\sk_{\algo{usr}}$, together with an interactive zero knowledge
proof of knowledge of its opening.

	\item The registration authority verifies that $\algo{uid}$ is fresh and
that the proof holds, and issues $\varphi$, a blind signature under $\rask$ on
the vector $(\algo{uid}, \sk_{\algo{usr}})$ committed in that message. The
registration authority does not learn $\sk_{\algo{usr}}$.

	\item $\algo{se}$ unblinds $\varphi$ and stores credential $\cred =
(\algo{uid}, \varphi)$, which satisfies $\DSverify(\rapk, (\algo{uid},
\sk_{\algo{usr}}), \varphi) = 1$.

	\item Pair the user's mobile device $\algo{dev}$ with $\algo{se}$, so
that only $\algo{dev}$ can send instructions to $\algo{se}$.

	\item Return $\cred$.

      \end{enumerate}

    \item[$T_0 \leftarrow \OPwithdraw(v_0)$ :] The user withdraws a token of
value $v_0$. Do the following:

      \begin{enumerate}

	\item The user's mobile device instructs their secure element to
initiate a withdrawal of value $v_0$.

	\item The secure element selects a random $s_0$, computes the pseudonym
$\com_0 \leftarrow \ComCommit(\ck, \algo{uid}_0, s_0)$, and returns $(s_0,
\com_0)$ to
the device.

	\item The device computes the request commitment $C$ to the vector
$(\algo{uid}_0, v_0, s_0)$, under the generator vector published as part of
$\pk_{\algo{cb}}$. This is a commitment to the three-message vector the central
bank will sign, and is a different object from the pseudonym $\com_0$, which
commits to $\algo{uid}_0$ alone under $\ck$.

	\item The device prepares a blind signature request on the committed
vector $(\algo{uid}_0, v_0, s_0)$, hiding $\algo{uid}_0$ and $s_0$ and
disclosing $v_0$. The request comprises $C$, the value $v_0$, the encryptions
$(\hat{\ciph}_0, \hat{\ciph}_2)$ of the two hidden messages under a freshly
generated key of the device, and a zero knowledge proof $\Pi$ that $C$ is
correctly formed and that those encryptions are consistent with it. The device
sends $\algo{req} = (C, v_0, \hat{\ciph}_0, \hat{\ciph}_2, \Pi)$ to the
intermediary.  Nothing in $\algo{req}$ reveals $\algo{uid}_0$ or $s_0$.

	\item The intermediary verifies $\Pi$, locks $v_0$ in the user's
account, and forwards $\algo{req}$ to the central bank.

	\item The central bank checks that $C$ is fresh, verifies $\Pi$, debits
the intermediary's reserves, and returns a blind signature response
$\algo{resp}$.

	\item The device unblinds $\algo{resp}$ to obtain $\psi$ with
$\DSverify(\pk_{\algo{cb}}, (\algo{uid}_0, v_0, s_0), \psi) = 1$, and stores
$(v_0, s_0, \com_0, \psi)$. The secure element stores the original token $\tau_0
= (v_0, \com_0)$.

	\item The device produces a proof of withdrawal $\Sigma$ with $\inst =
(\ck, \pk_{\algo{cb}}, v_0, \com_0)$, $\witness = (\algo{uid}_0, s_0, \psi)$,
and relation as follows: \begin{equation*} \DSverify(\pk_{\algo{cb}},
(\algo{uid}_0, v_0, s_0), \psi) = 1 \;\wedge\; \ComOpen(\ck, \com_0, s_0) =
\algo{uid}_0.  \end{equation*}

      \item Return $T_0 = (v_0, s_0, \com_0, \algo{hist}_0 = \Sigma)$.

    \end{enumerate}

  \item[$p_i \leftarrow \OPexecute(v, s_i, \com_i, \algo{hist}_i)$ :] The payer
pays the payee using token $(v, s_i, \com_i, \algo{hist}_i)$. Do the following:

    \begin{enumerate}

      \item The payee's device instructs their secure element to prepare to
receive a payment of value $v$.

      \item The payee's secure element selects a random $s_{i+1}$, computes
$\com_{i+1} \leftarrow \ComCommit(\ck, \algo{uid}_{i+1}, s_{i+1})$, and returns
$(s_{i+1}, \com_{i+1})$ to their device.

      \item The payee's device sends $\com_{i+1}$ to the payer's device.

      \item The payer's device notifies their secure element that a payment is
under way. The secure element locks token $\tau_i = (v, \com_i)$.

      \item The payer's device computes the audit ciphertext $\ciph_i \leftarrow
\PKEencrypt(\pk_\auditor, \algo{uid}_i; \rho_i)$ and a proof of correctness
$\Gamma_i$, with $\inst = (\ck, \pk_\auditor, \com_i, \ciph_i)$, $\witness =
(\algo{uid}_i, s_i, \rho_i)$, and relation as follows: \begin{equation*}
\ComOpen(\ck, \com_i, s_i) = \algo{uid}_i \;\wedge\; \ciph_i =
\PKEencrypt(\pk_\auditor, \algo{uid}_i; \rho_i).  \end{equation*}

    \item The payer's secure element and device jointly compute $\sigma_i$, a
signature of knowledge on the message $m_i = (v, \com_i, \com_{i+1})$ for the
relation as follows: \begin{equation*} \DSverify(\rapk, (\algo{uid}_i,
\sk_{\algo{usr}_i}), \varphi_i) = 1 \;\wedge\; \ComOpen(\ck, \com_i, s_i) =
\algo{uid}_i.  \end{equation*} Binding $\com_{i+1}$ into $m_i$ is what fixes
the
destination of the payment. This is a joint proof of knowledge in the sense
of~\Cref{sec:joint-pok-def}: the secure element takes the role of $\algo{Prv}_0$
and contributes knowledge of the payer's secret key, and the device takes the
role of $\algo{Prv}_1$ and contributes the remaining witness.

    \item The payer's device sends payment $p_i = (v, \com_i, \com_{i+1},
\ciph_i, \sigma_i, \Gamma_i, \algo{hist}_i)$ to the payee's device.

    \item The payee's device verifies $\algo{hist}_i$: if $i = 0$, it checks
$\Sigma$. Otherwise it verifies $\Sigma$ and every $(\Gamma_k, \sigma_k)_{k \in
[i]}$.

    \item If $\Gamma_i$ and $\sigma_i$ also verify, the payee's device accepts
and stores $T_{i+1} = (v, s_{i+1}, \com_{i+1}, \algo{hist}_{i+1} = (\com_i,
\ciph_i, \sigma_i, \Gamma_i) \, \| \, \algo{hist}_i)$, and their secure element
stores $\tau_{i+1} = (v, \com_{i+1})$.

    \item Upon acknowledgement, the payer's secure element deletes $\com_i$, and
their device returns $p_i$.

  \end{enumerate}

\item[$\{0,1\} \leftarrow \OPdeposit(\tau_n, \algo{hist}_n)$ :] The user
deposits token $\tau_n = (v, \com_n)$. Do the following:

  \begin{enumerate}

    \item The user instructs their device to deposit $\tau_n$. The device
informs the secure element to lock $\tau_n$.

    \item The device computes the audit ciphertext $\ciph_n \leftarrow
\PKEencrypt(\pk_\auditor, \algo{uid}_n; \rho_n)$ and proof $\Gamma_n$, exactly
as in $\OPexecute$.

    \item The device and secure element jointly compute $\sigma_n$, a signature
of knowledge on the message $m_n = (v, \com_n, \algo{dep})$ for the relation of
$\OPexecute$, where $\algo{dep}$ is a fixed marker distinct from every
pseudonym, so that a deposit signature cannot be replayed as a payment.

    \item The device sends deposit request $d_n = (\tau_n, \ciph_n, \sigma_n,
\Gamma_n, \algo{hist}_n)$ to the intermediary, which forwards it to the central
bank.

    \item The central bank verifies $\sigma_n$, $\Gamma_n$, and $\algo{hist}_n$.
If any check fails, return $0$.

    \item Otherwise, the central bank checks whether any prior deposit shares
the same origin $\com_0$. If so, it records double spending, notifies the
auditor, and returns $0$.

    \item Otherwise, the central bank credits the intermediary's account, which
in turn credits the user's account, and the secure element deletes $\tau_n$ upon
notification.

    \item Return $1$.

  \end{enumerate}

\item[$\algo{DS} \leftarrow \OPidentify(d_n, d_m)$ :] On input two conflicting
deposit requests $d_n$ and $d_m$ involving the same original token, the auditor
does:

  \begin{enumerate}

    \item Parse $d_n$ into the pseudonyms $(\com_k)_{k \in [n+1]}$ and
ciphertexts $(\ciph_k)_{k \in [n+1]}$ it carries, taking these from
$\algo{hist}_n$ together with the depositor's own $\com_n$ and $\ciph_n$, and
parse $d_m$ analogously into $(\com'_k)_{k \in [m+1]}$ and $(\ciph'_k)_{k \in
[m+1]}$.

    \item Find the largest index $j$ such that $\com_k = \com'_k$ for all $k
\leq j$. Either $\com_{j+1} \neq \com'_{j+1}$, in which case the owner at $j$
sent two different payments, or one of the two chains ends at $j$, in which case
the owner at $j$ deposited the token and also spent or deposited it a second
time. Both requests carry $\ciph_j$.

    \item Decrypt $\ciph_j$ using $\sk_\auditor$ to recover $\algo{uid}^*$, the
identifier of the owner of $\tau_j = (v, \com_j)$.

    \item Return $\algo{DS} = \{\algo{uid}^*\}$, extended by repeating the above
for each further pair of conflicting deposits of the same token.

  \end{enumerate}

\end{description}

\end{protocolbox}

\begin{theorem}\label{thm:offline-construction-security}
  %
  If the commitment scheme is binding, the public key encryption scheme is
IND-CPA secure, the blind signature scheme used by the central bank and the
signature scheme used by the registration authority are existentially
unforgeable, and the proofs and signatures of knowledge are simulation
extractable, with the withdrawal proof additionally online extractable, then
Protocol~\ref{prot:offline} securely realises ideal functionality
$\idfu{\algo{op}}$ against static corruption of devices and secure elements,
with the registration authority, the central bank, and the intermediaries honest
but curious.
%
\end{theorem}

\paragraph{Extensions.} The construction supports several extensions.

The construction extends to epoch based expiry, which is how it accommodates
revocation. The registration authority signs the triple $(\algo{uid},
\sk_{\algo{usr}}, \epoch)$ rather than the pair $(\algo{uid},
\sk_{\algo{usr}})$, and the central bank signs $(\algo{uid}_0, v, s_0, \epoch)$
at withdrawal, tying each token to the epoch in which it was issued. Both
schemes therefore sign vectors one longer than in
\Cref{sec:offline-instantiation}, and the withdrawal proof and the joint proof
gain one exponent each. Nothing else changes.  Payees enforce the policy,
rejecting payments carrying an epoch outside the window they accept, and the
width of that window can be tuned to the value of the payment. An epoch can
equally be counted in spending hops rather than in wall clock time, capping the
length of a chain that will be accepted at deposit. This is the form revocation
takes offline: a payee who has been offline since before a payer was revoked has
no verified real time information with which to reject that payer, so the
guarantee is expressed as a bound on how long a credential remains usable rather
than as an instruction delivered to payees. Epochs are disclosed, so they narrow
the anonymity set of a withdrawal to the users who withdrew the same
denomination in the same epoch, and a deployment should choose the epoch width
with that in mind.

Capping the value of an individual payment extends the construction in the same
direction and composes with epochs. Together the two bound the value a
compromised wallet can move before its tokens expire, making double spending
financially unattractive independently of whether the spender is subsequently
identified. Both are payee side checks and require no change to the protocol
messages.

The construction is parametric in the hardware that realises the secure element,
and admits several instantiations. Mobile TEEs protect application secrets from
a malicious device owner, but expose them to the application developer through
key management APIs, so key confidentiality does not by itself give integrity of
the logic that uses the key~\cite{trustzone,secureenclave}. Smart cards such as
Javacard and MIFARE protect the key and the logic together, in tamper resistant
isolated execution, which is what the single-use enforcement of
Protocol~\ref{prot:offline} requires~\cite{javacard,mifare}. We instantiate with
a smart card in~\Cref{sec:offline-evaluation} for this reason. A deployment
chooses among these according to the equipment cost it can bear, the system wide
value limits it sets, and the cost of attacking a single card. No hardware
assumption is unconditional, and the epoch and value caps above are what bound
the exposure when one fails.

The requirement that $\sk_{\algo{usr}}$ never leave the secure element is what
the joint proof of knowledge exists to satisfy, and it extends to a recovery
mechanism at the credential layer. A user who loses their device re-enrols with
the registration authority under a fresh pair $(\algo{uid}, \sk_{\algo{usr}})$,
restoring their ability to transact without migrating any state. The unspent
tokens held by the lost device are not recoverable, since no other party holds
the commitment randomness needed to spend them. This matches the behaviour of
cash, and of the deployed offline systems of~\Cref{sec:offline-related}.

We leave divisibility, aggregated deposits, and coin transparency to future
work.
